%% toc: P3 28 28.9
\chapter{Decidability by Elimination and Interpretation}
\label{sec:p03-w03-p03-022-decidability}

A first-order theory is decided when an algorithm answers, for every
sentence of its signature, whether the theory proves it. This chapter
gives two methods for obtaining such algorithms and two matching methods
for recognising when none exists. The first method eliminates quantifiers
effectively and then evaluates the remaining quantifier-free sentence. The
second translates the sentences about one structure into sentences about
another structure whose theory is already known to be decidable.

The inputs are a computable signature, a theory or a structure, and either
an explicit elimination procedure or an explicit interpretation. The
conclusions are decision procedures, completeness of the theory when the
quantifier-free sentences are all settled, and computable reductions that
carry undecidability in the opposite direction. The general notion of
quantifier elimination and the outline of the one-quantifier method are
in \cref{sec:qe-categoricity-01}, and the classifications that the
procedures below establish are listed in \cref{tab:p03-qe-catalog}. The
quotient interpretations used in the second half are those of
\cref{def:p03-imaginaries-interpretation}, with their truth lemma
\cref{thm:p03-imaginaries-quotient-truth}. Every proof below is given in
full in the annex; the proofs for fields also use algebraic facts proved in
\cref{sec:p03-w03-p03-020-fields} and the sign-table construction proved in
\cref{sec:p03-w03-p03-037-qe}.

\section{Effective elimination and decision}
\label{sec:p03-w03-p03-022-effective}

\begin{hbconvention}{def:p03-w03-p03-022-standing}{Standing conventions}
Signatures are first-order with equality and have computably presented
symbol sets, so formulas can be coded by natural numbers and the coding of
syntactic operations is computable. The quantifier-free formulas include
the two constants \(\top\) and \(\bot\). For a formula \(\varphi\),
\(\operatorname{fv}(\varphi)\) is its set of free variables. A
\emph{literal} is an atomic formula or the negation of one. A theory
\(T\) is a set of sentences; \(T\entails\sigma\) and \(T\derives\sigma\)
agree by soundness and completeness of the first-order calculus, and
\(T\) is \emph{decidable} when \(\{\sigma: T\entails\sigma\}\) is a
decidable set of codes. For a structure \(\mathcal A\), \(\operatorname{Th}(\mathcal A)\) is
the set of sentences true in \(\mathcal A\); a quantifier-elimination claim
about \(\mathcal A\) means one about \(\operatorname{Th}(\mathcal A)\).
\end{hbconvention}

\begin{hbdefinition}{def:p03-w03-p03-022-effective-qe}{Effective quantifier elimination}
A theory \(T\) \emph{eliminates quantifiers effectively} if there is an
algorithm that sends each formula \(\varphi\) to a quantifier-free formula
\(\psi\) with \(\operatorname{fv}(\psi)\subseteq\operatorname{fv}(\varphi)\) and
\(T\entails\forall\bar x\,(\varphi\leftrightarrow\psi)\), where \(\bar x\)
lists \(\operatorname{fv}(\varphi)\). An \emph{elimination step} for \(T\)
is an algorithm doing the same for every formula
\(\exists y\,\theta\) in which \(\theta\) is a finite conjunction of
literals.
\index{quantifier elimination!effective}
\end{hbdefinition}

The free-variable condition is load-bearing. Without it a sentence could be
``eliminated'' to a formula with new free variables, and nothing would be
decided.

\begin{hblemma}{thm:p03-w03-p03-022-one-step}{proved=annex:p03-w03-p03-022-one-step}{An elimination step suffices}
If \(T\) has an elimination step, then \(T\) eliminates quantifiers
effectively.
\end{hblemma}

\begin{hbtheorem}{thm:p03-w03-p03-022-decision}{proved=annex:p03-w03-p03-022-decision}{Decision from elimination}
Let \(T\) eliminate quantifiers effectively.
\begin{enumerate}
\item If the set of quantifier-free sentences \(\psi\) with
  \(T\entails\psi\) is decidable, then \(T\) is decidable.
\item If \(T\) is consistent and, for every atomic sentence \(\alpha\),
  either \(T\entails\alpha\) or \(T\entails\neg\alpha\), then \(T\) is
  complete. In particular this holds when the signature has no constant
  symbols and no nullary relation symbols, since there are then no atomic
  sentences.
\item If \(T\) satisfies the hypothesis of item 2, the decision of the
  atomic sentences is computable, and \(T\) is consistent, then \(T\) is
  complete and decidable, and \(T\) axiomatizes \(\operatorname{Th}(\mathcal A)\) for every
  model \(\mathcal A\) of \(T\).
\end{enumerate}
\end{hbtheorem}

The procedure has two separate inputs. Each can hold while the other
fails, and the theorem then gives nothing.

\begin{hbwarning}{def:p03-w03-p03-022-two-inputs}{Elimination alone does not decide}
Expand the signature \(\{0,S,+,\cdot\}\) of arithmetic by a relation symbol
\(R_\varphi\) of arity \(n\) for each formula \(\varphi(x_1,\ldots,x_n)\),
nullary when \(\varphi\) is a sentence, and interpret \(R_\varphi\) in
\(\N\) as the set that \(\varphi\) defines. In the expanded structure the
map \(\varphi\mapsto R_\varphi(x_1,\ldots,x_n)\) is a computable elimination
for the original formulas. A formula \(\chi\) of the expanded signature is
handled by first replacing each atom \(R_\varphi(\bar t)\) by the
corresponding instance of \(\varphi\), which yields an original formula
\(\chi'\) equivalent to \(\chi\) in the expanded structure, and then
returning \(R_{\chi'}\) applied to the free variables of \(\chi\).
Yet the quantifier-free sentence \(R_\sigma\) is true
exactly when \(\sigma\) is true in \(\N\), so deciding quantifier-free
sentences would decide \(\operatorname{Th}(\N;0,S,+,\cdot)\). That theory
contains \(\mathsf Q\), which is true in \(\N\), and is consistent; if it were
decidable it would be a consistent computably enumerable extension of
\(\mathsf Q\) with a decidable set of theorems, contrary to
\cref{thm:theoremhood-undecidable}. Item 1 of
\cref{thm:p03-w03-p03-022-decision} therefore genuinely needs its second
input.
\end{hbwarning}

\section{Elimination procedures}
\label{sec:p03-w03-p03-022-procedures}

Each procedure below is an elimination step in the sense of
\cref{def:p03-w03-p03-022-effective-qe}: it receives
\(\exists y\,\theta\) with \(\theta\) a conjunction of literals and returns a
quantifier-free equivalent. Two preparatory moves are shared. A conjunct of
\(\theta\) not containing \(y\) is moved outside the quantifier, using
\(\exists y\,(\chi\land\eta)\leftrightarrow(\chi\land\exists y\,\eta)\) when
\(y\notin\operatorname{fv}(\chi)\). And a quantifier distributes over a
disjunction, so a literal that the theory makes equivalent to a finite
disjunction of atoms may be expanded and the step applied to each disjunct.

\subsection{Dense linear orders}

\begin{hbmethod}{def:p03-w03-p03-022-dlo-step}{Elimination step for \(\mathsf{DLO}\)}
The signature is \(\{<\}\) and \(\mathsf{DLO}\) is as in
\cref{thm:dlo-qe}.
\begin{enumerate}
\item Replace \(\neg(u<v)\) by \(v<u\lor u=v\) and \(\neg(u=v)\) by
  \(u<v\lor v<u\); distribute, so that each disjunct is a conjunction of
  atoms.
\item Delete each conjunct \(y=y\). If a conjunct \(y<y\) occurs, return
  \(\bot\) for that disjunct.
\item If a conjunct \(y=x\) or \(x=y\) occurs with \(x\) distinct from
  \(y\), return the result of substituting \(x\) for \(y\) in the
  remaining conjuncts.
\item Otherwise the \(y\)-conjuncts are \(x_i<y\) for \(i\in L\) and
  \(y<z_j\) for \(j\in U\). Return
  \(\bigwedge_{i\in L}\bigwedge_{j\in U}x_i<z_j\), read as \(\top\) when
  \(L\) or \(U\) is empty.
\end{enumerate}
\end{hbmethod}

\begin{hbtheorem}{thm:p03-w03-p03-022-dlo-step}{proved=annex:p03-w03-p03-022-dlo-step}{Correctness and consequences for \(\mathsf{DLO}\)}
The method of \cref{def:p03-w03-p03-022-dlo-step} is an elimination step
for \(\mathsf{DLO}\). Hence \(\mathsf{DLO}\) eliminates quantifiers
effectively, is complete, and is decidable. In particular
\(\operatorname{Th}(\mathbb Q;<)\) is axiomatized by \(\mathsf{DLO}\).
\end{hbtheorem}

Item 4 uses every axiom: density when both bound sets are nonempty, the
absence of endpoints when one is empty, and nonemptiness when both are. If
any one is dropped, the equivalence of item 4 fails in some model; for
example, in the two-element order \(0<1\) the formula \(\exists y\,(x<y)\)
holds of \(0\) and fails of \(1\), while its proposed equivalent \(\top\)
holds of both.

\subsection{Successor with a named initial element}

\begin{hbdefinition}{def:p03-w03-p03-022-succ-theory}{The theory \(T_{\mathrm S}\)}
In the signature \(\{0,S\}\), with \(S\) unary, let \(S^0t=t\) and
\(S^{k+1}t=S(S^kt)\). The theory \(T_{\mathrm S}\) has the axioms
\[
 \forall x\,S x\ne 0,\qquad
 \forall x\forall y\,(Sx=Sy\to x=y),\qquad
 \forall x\,(x\ne0\to\exists y\,Sy=x),
\]
and, for every \(n\ge1\), \(\forall x\,S^nx\ne x\).
\end{hbdefinition}

\begin{hbmethod}{def:p03-w03-p03-022-succ-step}{Elimination step for \(T_{\mathrm S}\)}
Every term is \(S^kv\) with \(v\) a variable or \(0\), and every atom is
\(S^ku=S^lv\).
\begin{enumerate}
\item Replace a conjunct \(S^ky=S^ly\) by \(\top\) if \(k=l\) and by
  \(\bot\) otherwise.
\item Let \(K\) be the largest exponent of \(y\) in the remaining
  conjuncts. Replace each conjunct \(S^ky=S^lv\) or its negation, with
  \(v\) other than \(y\), by \(S^Ky=S^{l+K-k}v\) or its negation.
\item Replace \(S^Ky\) by a fresh variable \(z\), and add the conjuncts
  \(z\ne S^j0\) for \(j<K\). The formula is now
  \(\exists z\,\eta\) where every \(z\)-literal is \(z=t\) or \(z\ne t\)
  with \(t\) free of \(z\).
\item If some conjunct is \(z=t\), return the result of substituting
  \(t\) for \(z\) in the remaining conjuncts. Otherwise return \(\top\).
\end{enumerate}
\end{hbmethod}

\begin{hbtheorem}{thm:p03-w03-p03-022-succ}{proved=annex:p03-w03-p03-022-succ}{Elimination and decision for successor}
The method of \cref{def:p03-w03-p03-022-succ-step} is an elimination step
for \(T_{\mathrm S}\). Consequently \(T_{\mathrm S}\) eliminates quantifiers
effectively and is complete and decidable. It axiomatizes the theory
\(\operatorname{Th}(\N;0,S)\). The theory \(\operatorname{Th}(\N;S)\) in the
signature \(\{S\}\) is decidable.
\end{hbtheorem}

\begin{hbwarning}{def:p03-w03-p03-022-succ-no-zero}{Without the named element}
In the signature \(\{S\}\) alone, \(\operatorname{Th}(\N;S)\) does not
eliminate quantifiers. A quantifier-free formula in the single variable
\(x\) is a Boolean combination of atoms \(S^kx=S^lx\), each of which is
true at every element of \(\N\) or at none. So such a formula defines
\(\emptyset\) or \(\N\). The formula \(\exists y\,Sy=x\) defines
\(\N\setminus\{0\}\). The decision procedure of
\cref{thm:p03-w03-p03-022-succ} survives because a sentence in \(\{S\}\)
is also a sentence in \(\{0,S\}\) with the same truth value in \(\N\);
the elimination is carried out after the expansion.
\end{hbwarning}

\subsection{Atomless Boolean algebras}

\begin{hbdefinition}{def:p03-w03-p03-022-aba}{Atomless Boolean algebras}
The signature has constants \(0\) and \(1\), binary meet \(\sqcap\),
binary join \(\sqcup\) and unary complement \(-\). \(\mathsf{BA}\) is the theory of Boolean algebras
together with \(0\ne1\). Write \(u\le v\) for \(u\sqcap v=u\). The theory
\(\mathsf{ABA}\) adds
\[
 \forall x\,\bigl(x\ne0\to\exists y\,(y\le x\land y\ne0\land y\ne x)\bigr).
\]
For terms \(u,v\), put \(u\triangle v=(u\sqcap -v)\sqcup(-u\sqcap v)\).
\end{hbdefinition}

\begin{hbmethod}{def:p03-w03-p03-022-aba-step}{Elimination step for \(\mathsf{ABA}\)}
\begin{enumerate}
\item Replace each atom \(u=v\) by \(u\triangle v=0\). Each literal is now
  \(\tau=0\) or \(\tau\ne0\).
\item For each term \(\tau(\bar x,y)\) put \(\tau_1=\tau(\bar x,1)\) and
  \(\tau_0=\tau(\bar x,0)\). Combine the equations \(\tau^{(i)}=0\) into the
  pair \(a=0\) and \(b=0\) with \(a=\bigsqcup_i\tau^{(i)}_1\) and
  \(b=\bigsqcup_i\tau^{(i)}_0\), both \(0\) if there are no equations.
\item Replace each \(\sigma\ne0\) by
  \(\sigma_1\sqcap y\ne0\lor\sigma_0\sqcap-y\ne0\) and distribute. Each
  disjunct has \(y\)-literals \(a\sqcap y=0\), \(b\sqcap-y=0\),
  \(c_j\sqcap y\ne0\) for \(j\in J\), and \(d_k\sqcap -y\ne0\) for
  \(k\in K\).
\item Return
  \(a\sqcap b=0\land\bigwedge_{j\in J}c_j\sqcap-a\ne0
   \land\bigwedge_{k\in K}d_k\sqcap-b\ne0\).
\end{enumerate}
\end{hbmethod}

\begin{hbtheorem}{thm:p03-w03-p03-022-aba}{proved=annex:p03-w03-p03-022-aba}{Elimination and decision for atomless Boolean algebras}
The method of \cref{def:p03-w03-p03-022-aba-step} is an elimination step
for \(\mathsf{ABA}\). Consequently \(\mathsf{ABA}\) eliminates quantifiers
effectively, is complete and is decidable.
\end{hbtheorem}

The atomless axiom is used once, to split each atom of a finite partition
in item 4, and it cannot be removed. Let \(\mathcal P\) be the algebra of
subsets of \(\{1,2,3\}\). By induction on terms, every term \(t(x)\) in the
single variable \(x\) takes, as a function of \(x\), one of the four forms
\(0\), \(1\), \(x\), \(-x\) in every Boolean algebra, and whether an equation
between two of these forms holds at an element depends only on whether the
element is \(0\) or \(1\). Hence \(\{1\}\) and \(\{1,2\}\) satisfy the same
quantifier-free formulas of \(\mathcal P\), while
\(\exists y\,(y\le x\land y\ne0\land y\ne x)\) holds at \(\{1,2\}\) and
fails at the atom \(\{1\}\). So \(\operatorname{Th}(\mathcal P)\), a complete
extension of \(\mathsf{BA}\), does not eliminate quantifiers.

\subsection{Integers and naturals with addition and congruences}

\begin{hbdefinition}{def:p03-w03-p03-022-presburger}{The signatures \(\Lang_{\mathbb Z}\) and \(\Lang_{\N}\)}
\(\Lang_{\mathbb Z}\) has constants \(0,1\), binary \(+\), unary \(-\), binary
\(<\), and a unary relation symbol \(\mathrm D_n\) for each integer
\(n\ge1\). The structure \(\mathcal Z\) is \(\mathbb Z\) with the usual
interpretations and \(\mathrm D_n(m)\) iff \(n\) divides \(m\). The signature
\(\Lang_{\N}\) omits \(-\); the structure \(\mathcal N\) is \(\N\) with the
same interpretations. For an integer \(k\) and a term \(t\), \(kt\)
abbreviates \(t+\cdots+t\) (\(k\) summands) for \(k>0\), \(0\) for
\(k=0\), and \(-(|k|t)\) for \(k<0\); \(k\) alone abbreviates \(k1\).
\end{hbdefinition}

Every \(\Lang_{\mathbb Z}\)-term \(t\) has the same value in \(\mathcal Z\) as
a unique \emph{normal term} \(\sum_ik_ix_i+c\) with integer coefficients,
and the normal term is computed by collecting coefficients.

\begin{hbmethod}{def:p03-w03-p03-022-cooper-step}{Elimination step for \(\operatorname{Th}(\mathcal Z)\)}
\begin{enumerate}
\item Normalize terms. Replace \(s=t\) by \(0<t-s+1\land0<s-t+1\),
  \(s<t\) by \(0<t-s\), \(\neg(0<t)\) by \(0<1-t\), and
  \(\neg\mathrm D_n(t)\) by \(\bigvee_{r=1}^{n-1}\mathrm D_n(t+r)\) (by
  \(\bot\) when \(n=1\)); distribute. Each \(y\)-literal is now
  \(0<ky+t\) or \(\mathrm D_n(ky+t)\) with \(k\ne0\) and \(t\) free of
  \(y\).
\item Let \(\ell\) be the least common multiple of the \(|k|\). Multiply
  \(0<ky+t\) by \(\ell/|k|\). Replace \(\mathrm D_n(ky+t)\) by
  \(\mathrm D_{n\ell/|k|}(\ell y+(\ell/k)t)\). Every \(y\)-occurrence is now
  \(\ell y\) or \(-\ell y\).
\item Replace \(\ell y\) by a fresh \(z\) and add the conjunct
  \(\mathrm D_\ell(z)\). The \(z\)-literals are \(z<a_i\) (\(i\in I\)),
  \(b_j<z\) (\(j\in J\)) and \(\mathrm D_{m_h}(z+c_h)\) (\(h\in H\)), with
  \(a_i,b_j,c_h\) free of \(z\). Write \(\eta(z)\) for their conjunction.
\item Let \(M\) be the least common multiple of the \(m_h\). If \(J\) is
  nonempty, return \(\bigvee_{j\in J}\bigvee_{r=1}^M\eta(b_j+r)\). If
  \(J\) is empty, return
  \(\bigvee_{r=1}^M\bigwedge_{h\in H}\mathrm D_{m_h}(r+c_h)\).
\end{enumerate}
\end{hbmethod}

\begin{hbtheorem}{thm:p03-w03-p03-022-presburger}{proved=annex:p03-w03-p03-022-presburger}{Elimination and decision with congruence predicates}
The method of \cref{def:p03-w03-p03-022-cooper-step} is an elimination step
for \(\operatorname{Th}(\mathcal Z)\). Consequently \(\operatorname{Th}(\mathcal Z)\) eliminates
quantifiers effectively and is decidable. Moreover
\(\operatorname{Th}(\mathcal N)\) eliminates quantifiers effectively in
\(\Lang_{\N}\) and is decidable, and so is
\(\operatorname{Th}(\N;0,1,+,<)\).
\end{hbtheorem}

\begin{hbexample}{def:p03-w03-p03-022-parity}{Eliminating a parity quantifier}
Take \(\exists y\,(y+y=x)\). Item 1 gives
\(\exists y\,(0<x-2y+1\land0<2y-x+1)\). Here \(\ell=2\); item 3 gives
\(\exists z\,(z<x+1\land x-1<z\land\mathrm D_2(z))\), so
\(J\) contains the single bound \(b=x-1\), \(H\) contains \(\mathrm D_2(z+0)\),
and \(M=2\). Item 4 returns the disjunction over \(r=1,2\) of
\(x-1+r<x+1\land x-1<x-1+r\land\mathrm D_2(x-1+r)\). Evaluating the
order literals, the disjunct \(r=1\) reduces to \(\mathrm D_2(x)\) and the
disjunct \(r=2\) contains \(x+1<x+1\), which is false. The result is
equivalent to \(\mathrm D_2(x)\).
\end{hbexample}

\begin{hblemma}{thm:p03-w03-p03-022-presburger-boundary}{proved=annex:p03-w03-p03-022-presburger-boundary}{The congruence predicates are needed}
In the signature \(\{0,1,+,<\}\), no quantifier-free formula \(\psi(x)\)
satisfies \(\mathcal N\models\psi(m)\) exactly for the even \(m\). Hence
\(\operatorname{Th}(\N;0,1,+,<)\) does not eliminate quantifiers in that
signature, although it is decidable.
\end{hblemma}

\subsection{Algebraically closed fields}
\label{sec:p03-w03-p03-022-acf}

The language is the ring language \(\Lang_{\mathrm r}\) of
\cref{sec:p03-w03-p03-020-fields}, and \(\mathsf{ACF}_p\) is the theory
defined there. Write \(\mathsf{ACF}\) for the field axioms together with
the root axioms \(\forall a_0\cdots\forall a_{d-1}\,\exists y\,
(y^d+a_{d-1}y^{d-1}+\cdots+a_0=0)\) for all \(d\ge1\), with no axiom on the
characteristic. Every \(\Lang_{\mathrm r}\)-term \(t\) denotes, in every
commutative ring, the value of an integer polynomial \(P_t\), computed from
\(t\) by recursion on \(t\); conversely, a term denoting a given integer
polynomial is computed from it. For an integer polynomial \(c\), the
formula \(c=0\) is the atom \(t=0\) with \(P_t=c\).

\begin{hbconvention}{def:p03-w03-p03-022-parametric}{Parametric polynomials over a field}
Parametric polynomials in a variable \(y\) with parameter variables
\(\bar x\), their formal degree, formal leading coefficient and truncations
\(F_{\le j}\) are as in \cref{def:p03-w03-p03-037-parametric}, with
coefficients in \(\mathbb Z[\bar x]\). \emph{Padding} \(F\) to a formal
degree \(N\) at least its own appends zero coefficients. For a field \(K\)
and \(\bar a\in K^m\), \(F(\bar a)\in K[y]\) is obtained by evaluating the
coefficients at \(\bar a\) through the ring homomorphism
\(\mathbb Z\to K\). The pseudo-remainder \(\operatorname{prem}(H,F)\) is
that of \cref{def:p03-w03-p03-037-prem}, computed by the recursion that
proves \cref{thm:p03-w03-p03-037-prem}; that recursion also covers the
case in which \(H\) has formal degree \(q-1\), where it returns \(H\)
itself. So when \(F\) has formal degree
\(q\ge1\) and formal leading coefficient \(b\), and \(H\) has formal degree
\(N\ge q-1\), it has formal degree \(q-1\) and satisfies
\(b^{N-q+1}H=AF+\operatorname{prem}(H,F)\) in \(\mathbb Z[\bar x][y]\) for
some \(A\) (\cref{thm:p03-w03-p03-037-prem}). This is an identity between
integer polynomials, so it survives evaluation in every field.
\end{hbconvention}

\begin{hbmethod}{def:p03-w03-p03-022-acf-step}{Elimination step for \(\mathsf{ACF}\)}
\emph{Input.} \(\exists y\,\theta\) with \(\theta\) a conjunction of
\(\Lang_{\mathrm r}\)-literals; \(\bar x\) lists the other variables of
\(\theta\).

\emph{Preparation.} For each literal \(t=s\) or \(\neg\,t=s\), compute
\(P_t-P_s\) and read it as a parametric polynomial in \(y\) whose formal
degree is its degree in \(y\) (formal degree \(0\) for the zero
polynomial). Let \(F_1,\ldots,F_k\) come from the equations and
\(G_1,\ldots,G_l\) from the negated equations, and let \(G\) be the product
\(G_1\cdots G_l\) with the sum of their formal degrees (\(G=1\), of formal
degree \(0\), when \(l=0\)). Return \(\operatorname{Elim}(F_1,\ldots,F_k;G)\).

\emph{Recursion.} \(\operatorname{Elim}(F_1,\ldots,F_k;G)\) is given by the
first applicable clause.
\begin{enumerate}
\item If \(k=0\), return \(\bigvee_c\neg\,c=0\), where \(c\) ranges over
  the coefficients of \(G\).
\item If some \(F_i\) has formal degree \(0\), say \(F_i=c\) with \(i\)
  least, return
  \(c=0\land\operatorname{Elim}(F_1,\ldots,F_{i-1},F_{i+1},\ldots,F_k;G)\).
\item If \(k=1\), let \(F=F_1\) have formal degree \(q\ge1\) and formal
  leading coefficient \(b\), let \(g\) be the formal degree of \(G\), let
  \(H\) be \(G^q\) padded to formal degree \(\max(qg,q-1)\), and let
  \(R=\operatorname{prem}(H,F)\). Return
  \[
  \bigl(b=0\land\operatorname{Elim}(F_{\le q-1};G)\bigr)\lor
  \Bigl(\neg\,b=0\land\bigvee_r\neg\,r=0\Bigr),
  \]
  where \(r\) ranges over the coefficients of \(R\).
\item If \(k\ge2\), let \(i\) be least such that \(F_i\) has the least
  formal degree \(q\), and let \(b\) be its formal leading coefficient.
  Return
  \begin{multline*}
  \bigl(b=0\land\operatorname{Elim}(F_1,\ldots,(F_i)_{\le q-1},\ldots,F_k;G)\bigr)\\
  \lor\bigl(\neg\,b=0\land\operatorname{Elim}(F'_1,\ldots,F'_k;G)\bigr),
  \end{multline*}
  with \(F'_i=F_i\) and \(F'_j=\operatorname{prem}(F_j,F_i)\) for \(j\ne i\).
\end{enumerate}
Side conditions: \(\operatorname{prem}(F_j,F_i)\) is defined in item 4
because \(F_j\) has formal degree at least \(q\). The pair (sum of the
formal degrees of the \(F_i\), number \(k\)) decreases lexicographically
at every recursive call, so the recursion stops. Item 3 tests whether
\(F(\bar a)\) divides \(G(\bar a)^q\).
\end{hbmethod}

\begin{hbtheorem}{thm:p03-w03-p03-022-acf-step}{proved=annex:p03-w03-p03-022-acf-step}{Correctness of the step}
For every model \(K\) of \(\mathsf{ACF}\), every list
\(F_1,\ldots,F_k,G\) of parametric polynomials with parameters \(\bar x\)
and every \(\bar a\in K^m\), the formula
\(\operatorname{Elim}(F_1,\ldots,F_k;G)\) holds in \(K\) at \(\bar a\)
if and only if some \(y\in K\) satisfies
\[
F_1(\bar a)(y)=\cdots=F_k(\bar a)(y)=0
\quad\text{and}\quad G(\bar a)(y)\ne0 .
\]
Consequently the method of \cref{def:p03-w03-p03-022-acf-step} is an
elimination step for \(\mathsf{ACF}\) and for every \(\mathsf{ACF}_p\).
\end{hbtheorem}

\begin{hbtheorem}{thm:p03-w03-p03-022-acf}{proved=annex:p03-w03-p03-022-acf}{Decision for algebraically closed fields}
Assume the Axiom of Choice. For \(p=0\) and every prime \(p\),
\(\mathsf{ACF}_p\) eliminates quantifiers effectively in
\(\Lang_{\mathrm r}\), is complete and decidable, and axiomatizes
\(\operatorname{Th}(K)\) for every algebraically closed field \(K\) of
characteristic \(p\). The theory \(\mathsf{ACF}\) eliminates quantifiers
effectively and is decidable, but it is not complete.
\end{hbtheorem}

The choice principle enters only through
\cref{thm:p03-w03-p03-020-closure-exists}, which supplies a model of each
\(\mathsf{ACF}_p\); the procedure and its correctness do not use it. The
theorem makes the elimination of \cref{thm:p03-w03-p03-020-acf-qe}
effective; that proof runs through the one-point test and yields no
algorithm. For \(\mathsf{ACF}\), item 1 of
\cref{thm:p03-w03-p03-022-decision} applies while item 2 fails, since the
atomic sentences \(\chi_p\) are left undecided.

\begin{hbexample}{def:p03-w03-p03-022-acf-example}{Two runs of the step}
Take \(\exists y\,(y\cdot y=x\land\neg\,y=0)\). Preparation gives
\(F=y^2-x\), with coefficients \((-x,0,1)\) and formal degree \(2\), and
\(G=y\) of formal degree \(1\). Item 3 applies with \(q=2\) and \(b=1\).
Here \(H=y^2\) has formal degree \(2\), and the recursion of
\cref{thm:p03-w03-p03-037-prem} performs one reduction,
\(1\cdot H_{\le1}-1\cdot F_{\le1}=x+0\cdot y\), of formal degree \(1\); so
\(R=x+0\cdot y\). The output is
\((1=0\land\operatorname{Elim}(F_{\le1};G))\lor(\neg\,1=0\land(\neg\,x=0\lor\neg\,0=0))\).
In every field the first disjunct and \(\neg\,0=0\) are false, so the
output is equivalent to \(\neg\,x=0\): a nonzero element has a nonzero
square root.

Take \(\exists y\,(y\cdot y=1\land y=x)\), with \(F_1=y^2-1\),
\(F_2=y-x\) and \(G=1\). Item 4 picks \(F_2\), with \(q=1\) and \(b=1\).
Two reductions give \(\operatorname{prem}(F_1,F_2)=x^2-1\), of formal degree
\(0\). Item 2 then outputs \(x\cdot x-1=0\) conjoined with
\(\operatorname{Elim}(F_2;1)\). Item 3 applied to \(F_2\) and \(G=1\) gives
\(H=R=1\) and the output
\((1=0\land\operatorname{Elim}(-x;1))\lor(\neg\,1=0\land\neg\,1=0)\), which
is true in every field. The whole output is therefore equivalent, in every
field, to \(x\cdot x-1=0\).
\end{hbexample}

\begin{hbwarning}{def:p03-w03-p03-022-acf-boundary}{Where algebraic closure is used}
The correctness proof uses the root axioms twice: to split \(F(\bar a)\)
into linear factors in item 3, and, through the infiniteness of
algebraically closed fields, to find a non-root of \(G(\bar a)\) in item 1.
Neither use can be dropped. On \(\exists y\,(y\cdot y+1=0)\) item 3 returns
\((1=0\land\operatorname{Elim}(F_{\le1};1))\lor(\neg\,1=0\land(\neg\,1=0\lor\neg\,0=0))\),
which is true in every field, while the input fails in the ordered field
\(\mathbb R\). On \(\exists y\,\neg\,(y^p-y=0)\), with \(y^p\) the
\(p\)-fold product and \(p\) prime, item 1 returns a disjunction with the
disjunct \(\neg\,1=0\), true in every field, while the input fails in the
field with \(p\) elements, where every element is a root of \(y^p-y\):
\(0\) trivially, and each nonzero element because its multiplicative order
divides \(p-1\), the order of the group of nonzero elements. So the method is not an elimination step for the
theory of fields, in agreement with \cref{def:p03-w03-p03-020-quadratic},
which shows by preservation that no elimination exists there.
\end{hbwarning}

\subsection{Real closed fields}
\label{sec:p03-w03-p03-022-rcf}

For \(\mathsf{RCF}\) in the ordered-ring language \(\Lang_{\mathrm{or}}\) of
\cref{def:p03-w03-p03-037-rcf}, the elimination of
\cref{thm:p03-w03-p03-037-qe} is built from the one-step method of
\cref{def:p03-w03-p03-037-elimination-step}. A quantifier \(\exists x\) over
a conjunction of sign conditions is replaced by the disjunction of the
formulas \(\theta_{\mathcal P,t}\) of \cref{thm:p03-w03-p03-037-uniform-sign}
over the sign tables \(t\) that contain the required sign vector. Each
\(\theta_{\mathcal P,t}\) is defined by recursion on the list
\(\mathcal P\), from formal derivatives, coefficients and pseudo-remainders,
through a map \(\Phi\) on finite words. The next theorem records that every
ingredient is computable and draws the consequences.

\begin{hbtheorem}{thm:p03-w03-p03-022-rcf}{proved=annex:p03-w03-p03-022-rcf}{Effective elimination and decision for real closed fields}
The construction of \cref{thm:p03-w03-p03-037-qe} is an algorithm: it
sends each \(\Lang_{\mathrm{or}}\)-formula \(\varphi\) to a quantifier-free
\(\psi\) with \(\operatorname{fv}(\psi)\subseteq\operatorname{fv}(\varphi)\)
and \(\mathsf{RCF}\entails\forall\bar x\,(\varphi\leftrightarrow\psi)\).
Consequently \(\mathsf{RCF}\) eliminates quantifiers effectively, is
complete and decidable, and axiomatizes \(\operatorname{Th}(K)\) for every
real closed field \(K\). In particular
\(\operatorname{Th}(\mathbb R;0,1,+,-,\cdot,<)\) is decidable, and so is
\(\operatorname{Th}(\mathbb R;0,1,+,-,\cdot)\).
\end{hbtheorem}

\begin{hbexample}{def:p03-w03-p03-022-rcf-example}{The sign-table step on a linear equation}
Take \(\exists x\,(a\cdot x+b=0)\), so \(\mathcal P=(P)\) with
\(P=ax+b\) of formal degree \(1\), and the required sign vector is \((0)\).
Where \(a=0\), the list \(\mathcal P^\flat=(b)\) has one-letter sign tables,
and \(\theta_{\mathcal P^\flat,((0))}\) is \(b=0\). Where \(a\ne0\), the list
\(\mathcal Q\) of \cref{def:p03-w03-p03-037-elimination-step} is
\((P',a,a)\) with \(P'=a\): the formal derivative, the formal leading
coefficient and the one coefficient of \(P'\); there are no
pseudo-remainders, since \(P'\) has formal degree \(0\). Its sign table is
the single letter \((s,s,s)\) with \(s=\operatorname{sgn}a\). The map
\(\Phi\) finds no root of \(P'\), gives \(P\) the signs \(-s\) and \(s\) at
\(-\infty\) and \(+\infty\), and the fourth case of
\cref{thm:p03-w03-p03-037-interval} splits the line at one root; so
\(\Phi(((s,s,s)))=((-s),(0),(s))\). Every other one-letter word \(u\) has a
zero coordinate or two coordinates of opposite sign, so
\(\neg\,a=0\land\theta_{\mathcal Q,u}\) is false in every ordered field.
Deleting these false disjuncts and repeated conjuncts, the output is
\((a=0\land b=0)\lor(\neg\,a=0\land a>0)\lor(\neg\,a=0\land a<0)\), that
is, \(a\ne0\lor b=0\). For the sentence \(\exists x\,((1+1)\cdot x+1=0)\)
the same run has the integers \(2\) and \(1\) in place of \(a\) and \(b\),
and evaluating its integer atoms gives the value true; so \(\mathsf{RCF}\)
proves the sentence.
\end{hbexample}

The procedures of this subsection and the previous one carry no time bound.
The complexity entries of \cref{tab:p03-qe-catalog} are separate results
and are not derived from them.

\section{Interpretations and decidability transfer}
\label{sec:p03-w03-p03-022-interpretations}

\begin{hbdefinition}{def:p03-w03-p03-022-interpreted-structure}{A structure interpreted in another}
Let \(\mathcal A\) be a \(K\)-structure and \(\mathcal B\) an \(L\)-structure,
both one-sorted. An \emph{interpretation of \(\mathcal A\) in \(\mathcal B\)}
is an interpretation \(I\) in the sense of
\cref{def:p03-imaginaries-interpretation}, taken with source theory
\(\operatorname{Th}(\mathcal B)\) and one target sort, together with an isomorphism
\(I(\mathcal B)\cong\mathcal A\). It is \emph{computable} when \(K\) is
computably presented and the assignment of the domain, equality, relation
and graph formulas to the symbols of \(K\) is computable; when \(K\) is
finite every interpretation is computable.
\index{interpretation!decidability transfer}
\end{hbdefinition}

The parameter-free requirement is part of the definition. Parameters are
allowed only after naming them by constants in \(L\), as in
\cref{def:p03-imaginaries-standing}; the transfer below then concerns the
theory of the expanded \(\mathcal B\).

\begin{hbtheorem}{thm:p03-w03-p03-022-transfer}{proved=annex:p03-w03-p03-022-transfer}{Transfer of decidability and undecidability}
Let \(I\) be a computable interpretation of \(\mathcal A\) in \(\mathcal B\).
The map \(\sigma\mapsto\sigma^I\) of \cref{def:p03-imaginaries-translation}
is computable on \(K\)-sentences and satisfies
\[
  \mathcal A\models\sigma\quad\Longleftrightarrow\quad
  \mathcal B\models\sigma^I .
\]
Hence \(\operatorname{Th}(\mathcal A)\) is many-one reducible to
\(\operatorname{Th}(\mathcal B)\). If \(\operatorname{Th}(\mathcal B)\) is decidable,
so is \(\operatorname{Th}(\mathcal A)\); if \(\operatorname{Th}(\mathcal A)\) is
undecidable, so is \(\operatorname{Th}(\mathcal B)\).
\end{hbtheorem}

\begin{hbexample}{def:p03-w03-p03-022-nat-in-int}{The naturals inside the integers}
Interpret \(\mathcal N\) in \(\mathcal Z\) in dimension one with domain
\(\neg\,x<0\), equality \(x=x'\), and the relation and graph formulas of
\(\Lang_{\N}\) read literally in \(\Lang_{\mathbb Z}\). The identity map is
the required isomorphism. \Cref{thm:p03-w03-p03-022-transfer} yields the
decidability of \(\operatorname{Th}(\mathcal N)\) from that of
\(\operatorname{Th}(\mathcal Z)\). The relativization of quantifiers to
\(\neg\,x<0\) in the annex proof of \cref{thm:p03-w03-p03-022-presburger}
is this translation with its term-value clauses removed, which is possible
because the domain is closed under the operations.
\end{hbexample}

\begin{hbexample}{def:p03-w03-p03-022-int-in-nat}{The integers as differences of naturals}
Interpret \((\mathbb Z;0,1,+,-,<)\) in \((\N;0,1,+,<)\) in dimension two.
The domain is all pairs, the equality formula is
\(x_1+x_2'=x_2+x_1'\), and the formulas for \(0\), \(1\), \(+\), \(-\)
and \(<\) are \(z_1=z_2\), \(z_1=z_2+1\),
\(z_1+x_2+y_2=z_2+x_1+y_1\), \(z_1+x_1=z_2+x_2\) and
\(x_1+y_2<y_1+x_2\). The map \((p,q)\mapsto p-q\) is the isomorphism. Here
the equality formula is not literal equality: distinct pairs represent the
same integer, which is the reason quotient interpretations are needed.
\end{hbexample}

\begin{hbwarning}{def:p03-w03-p03-022-transfer-boundary}{What the transfer does not give}
The reduction runs from \(\mathcal A\) to \(\mathcal B\) only.
\((\N;S)\) is interpreted in \((\N;0,S,+,\cdot)\) by the identity
translation, its theory is decidable by \cref{thm:p03-w03-p03-022-succ},
and the theory of the larger structure is undecidable by
\cref{thm:theoremhood-undecidable}, by the argument of
\cref{def:p03-w03-p03-022-two-inputs}; nothing passes from the interpreted
structure back to the interpreting one. The theorem also concerns complete
theories of structures. For an interpretation of one incomplete theory in
another, plain undecidability does not transfer and the stronger
hypothesis of essential undecidability is required, as in
\cref{thm:p05-interpretation-transfer} \cite{TMR1953}.
\end{hbwarning}
